\setcounter{ex}{0}
\setcounter{paragraph}{0}
\PhanI
%\loai{Tính vận tốc cực đại và gia tốc cực đại}
\begin{ex}%[3P1B2-1][Loai1]
Một chất điểm dao động điều hòa với biên độ 10 cm và tần số góc 2 rad/s. Tốc độ cực đại của chất điểm là
\choice
{10 cm/s}
{40 cm/s}
{5 cm/s}
{\True 20 cm/s}
\loigiai{
Tốc độ cực đại của chất điểm là: $v_{\max} = \omega A = 2.10 = 20$ cm/s.
}
%\haidong{4}
\end{ex} \haidong{15}
\begin{ex}
Một chất điểm dao động điều hoà với tần số 4 Hz và biên độ dao động 10 cm. Độ lớn gia tốc cực đại của chất điểm bằng
\choice
{$2,5\ m/s^2$}
{$25\ m/s^2$}
{\True $63,2\ m/s^2$}
{$6,32\ m/s^2$}
\loigiai{
$a_{\max}=\omega^2A=(2\pi f)^2A=63,2\ m/s^2$.
}
%\haidong{4}
\end{ex} \haidong{15}
\begin{ex}%[3P1K2-1][Loai1]
Một vật nhỏ dao động điều hòa với biên độ 10 cm và vận tốc có độ lớn cực đại là 100 cm/s. Gia tốc cực đại của vật nhỏ là
\choice
{\True 10 $m/s^2$}
{1 $m/s^2$}
{1000 $m/s^2$}
{100 $ cm/s^2$}
\loigiai{
Ta có: $v_{\max} = \omega A \Rightarrow \omega  = 10\ rad/s \Rightarrow a_{\max} = \omega^2A = 10^2.10 = 1000\ cm/s^2 = 10\ m/s^2$.
}
%\haidong{4}
\end{ex} \haidong{15}
%\loai{Từ vận tốc cực đại hoặc gia tốc cực đại tìm các đại lượng}
\begin{ex}%[3P1B2-1][Loai1]
Một vật nhỏ dao động điều hòa với biên độ 5 cm và vận tốc có độ lớn cực đại là $10\pi$  cm/s. Chu kì dao động của vật nhỏ là
\choice
{4 s}
{2 s}
{\True 1 s}
{3 s}
\loigiai{
Ta có: $v_{\max} = \omega A \Rightarrow \omega  = \dfrac{v_{\max}}{A} = 2\pi  \ rad/s \Rightarrow T = 1$ s.
}
%\haidong{4}
\end{ex} \haidong{15}
%\loai{Từ vận tốc cực đại và gia tốc cực đại tìm các đại lượng}
\begin{ex}
Vận tốc của một vật dao động điều hoà khi đi qua vị trí cân bằng là 1 cm/s và gia tốc của vật khi ở vị trí biên là $1,57\ cm/s^2$. Chu kì dao động của vật là
\choice
{3,24 s}
{6,28 s}
{\True 4 s}
{2 s}
\loigiai{
\begin{itemize}
\item $\omega=\dfrac{a_{\max}}{v_{\max}}\dfrac{\pi}{T}\Rightarrow T=4\ s$.
\end{itemize}
}
%\haidong{4}
\end{ex} \haidong{15}
\begin{ex}%[3P1K2-1][Loai1]
Một vật nhỏ dao động điều hòa với gia tốc cực đại bằng 86,4 $m/s^2$, vận tốc cực đại bằng 2,16 m/s. Chiều dài quỹ đạo dao động của vật bằng
\choice
{5,4 cm}
{\True 10,8 cm}
{6,2 cm}
{12,4 cm}
\loigiai{
\begin{itemize}
\item Ta có: $\omega =\dfrac{a_{\max}}{{{v}_{\max }}}= 40\ rad/s \Rightarrow A = \dfrac{{{v}_{\max }}}{\omega } = 0,054\ m = 5,4\ cm$.
\item Quỹ đạo dao động là: $L = 2A = 10,8$ cm.
\end{itemize}
}

%\haidong{4}
\end{ex} \haidong{15}
\begin{ex}%[3P1K2-1][Loai1]
Một vật dao động điều hòa dọc theo trục Ox. Vận tốc cực đại của vật là $v_{\max} = 8{\pi}$  cm/s và gia tốc cực đại $a_{\max} = 16{\pi ^2}\  cm/s^2$. Trong thời gian một chu kì dao động vật đi được quãng đường là
\choice
{8 cm}
{12 cm}
{20 cm}
{\True 16 cm}
\loigiai{
\begin{itemize}
\item Ta có: $v_{\max} = 8\pi\  cm/s = \omega A;\ a_{\max} = 16\pi ^2 = \omega^2A \Rightarrow \omega  = 2\pi \Rightarrow A = 4$ cm.
\item Quãng đường vật đi trong 1 chu kì là $4A = 16$ cm.
\end{itemize}
}
%\haidong{4}
\end{ex} \haidong{15}

%\loai{Tìm vận tốc và gia tốc tại thời điểm t}
\begin{ex} %[3P1B2-1][Loai1] 
Một vật dao động điều hòa có phương trình: $x = 2\cos\left(2\pi t - \dfrac{\pi}{6}\right)$ (cm, s). Lấy $\pi^2 = 10$. Gia tốc của vật lúc $t = 0,25$ s là
\choice 
{$40\ cm/s^2$}
{\True $-40\ cm/s^2$}
{$\pm 40\ cm/s^2$}
{$-4\pi\ cm/s^2$}
\loigiai{
\begin{itemize}
\item Ta có: $a_{0,25} = -\omega^2.x_{0,25} = -40\ cm/s^2$.
\end{itemize}
}
%\haidong{4}
\end{ex} \haidong{15}
%\loai{Tổng hợp}

\begin{ex}%[3P1B2-1][Loai1]
Một vật dao động điều hoà trên trục Ox với phương trình  $x=5\sin \left( 4t+\dfrac{\pi}{6} \right)$ cm. Phương trình vận tốc là
\choice
{\True $v=20\cos \left(4t+\dfrac{\pi}{6}\right)\ cm/s$}
{$v=20\cos \left(4t+\dfrac{5\pi}{6}\right)\ cm/s$}
{$v=5\cos \left(4t+\dfrac{\pi}{3}\right)\ cm/s$}
{$v=20\cos \left(4t-\dfrac{\pi}{3}\right)\ cm/s$}
\loigiai{ 
\begin{itemize}
\item Phương trình dao động: $x=5\sin \left( 4t+\dfrac{\pi}{6} \right)=5\cos \left( 4t-\dfrac{\pi}{3} \right)$.
\item Phương trình vận tốc là: $v=\omega.A.\cos\left(\omega t+\varphi+\dfrac{\pi}{2}\right)=20\cos \left(4t+\dfrac{\pi}{6}\right)\ cm/s$.
\end{itemize}
}
%\haidong{3}
\end{ex} \haidong{15}
\begin{ex}%book
Phương trình vận tốc của một vật dao động là: $v=120\cos 20t\ (cm/s)$. Vào thời điểm $t=\dfrac{T}{6}$, vật có li độ là
\choice
{$3\ cm$}
{$-3\ cm$}
{\True $3\sqrt 3\ cm$}
{$-3\sqrt 3\ cm$}
\loigiai{
\begin{itemize}
\item $x=60\cos(20t-\dfrac{\pi}{2})=-3\sqrt 3\ cm$.
\end{itemize}
}
%\haidong{3}
\end{ex} \haidong{15}
\begin{ex}%book%[3P1K2-1][Loai1]
Một chất điểm dao động điều hòa với phương trình $x = 6\cos\pi t$ (x tính bằng cm, t tính bằng s). Phát biểu nào sau đây đúng?
\choice
{Chu kì của dao động là $0,5\ s$}
{\True Tốc độ cực đại của chất điểm là $18,8\ cm/s$}
{Gia tốc của chất điểm có độ lớn cực đại là $113\ cm/s^2$}
{Tần số dao động là $2\ Hz$}
\loigiai{
\begin{itemize}
\item Chu kì: $T = \dfrac{2\pi}{\omega}  = 2\ s \Rightarrow$ A sai.
\item Tần số: $f = \dfrac{1}{T} = 0,5$ Hz $\Rightarrow$ D sai.
\item Tốc độ cực đại: $v_{\max} = \omega A = 6\pi  \approx 18,8$ cm/s $\Rightarrow$ B đúng.
\item Gia tốc cực đại: $a_{\max} = \omega^2A = 6\pi ^2 \approx 59,16\ cm/s^2 \Rightarrow$ C đúng.
\end{itemize}
}
%\haidong{4}
\end{ex} \haidong{15}
\begin{ex}
Một vật dao động điều hòa với chu kì là $2\ s$ và biên độ dao động là $10\ cm$. Độ lớn vận tốc cực đại của vật là
\choice
{$20\pi\ cm/s$}
{$10\pi\ m/s$}
{\True $10\pi\ cm/s$}
{$20\pi\ m/s$}
\end{ex}
\begin{ex}
Một con lắc dao động điều hoà có phương trình $x = 8\cos\left(4\pi t + \dfrac{\pi}{3}\right)\ (cm)$.  
Vận tốc cực đại của con lắc có giá trị  là
\choice
{ $25,12\ cm/s$ }
{\True $100,48\ cm/s$ }
{ $50,24\ cm/s$ }
{ $12,56\ cm/s$ }
\loigiai{
}
\end{ex}
\begin{ex}
Một con lắc dao động điều hoà có phương trình $x = 8\cos\left(4\pi t + \dfrac{\pi}{3}\right)\ (cm)$.  
Gia tốc cực đại của con lắc có giá trị  là
\choice
{$157,7\ cm/s^{2}$ }
{$628,0\ cm/s^{2}$ }
{\True $1262\ cm/s^{2}$ }
{$314,0\ cm/s^{2}$ }
\loigiai{
}
\end{ex}
\begin{ex}
Quan sát đồ thị vận tốc – thời gian của một vật dao động điều hòa. Vận tốc cực đại của vật là\\
\centerline{\begin{tikzpicture}[very thick,>=stealth']
\pgfplotsset{width=9cm,height=4.5cm,compat=1.4}
\begin{axis}[
axis lines=middle, xmin=0, xmax=3.7*pi, xstep=1,
ymin=-4.5, ymax=5, ystep=0.1,
grid=major,%Vẽ chế độ lưới theo trục tọa độ Oxy,
x grid style={dashed, black},
axis line style={very thick,Mapcolor}, % <--- dòng thêm vào
xlabel=$t~(ms)$,
xlabel style={above},
xtick ={2,4,6,8,10},
xticklabels={},
y grid style={dashed, black},
extra y ticks={0},
extra y tick label={0},
ylabel=$v~(cm/s)$,
ylabel style={above},
ytick={-4, -2, 0, 2, 4},
yticklabels={$-4{,}2$, $-2{,}1$, $0$, $2{,}1$, $4{,}2$},]
\addplot[line width=1.2pt, domain=0:3.5*pi, samples=400, Mapcolor]{4*cos(deg(0.5*pi*x-pi/2))};
\addplot[Mapcolor,mark=*,mark options={fill=white,mark size=1.5pt}, nodes near coords={0,33}]
coordinates {(2, -4)};
\addplot[Mapcolor,mark=*,mark options={fill=white,mark size=1.5pt}, nodes near coords={0,66}]
coordinates {(4, -4)};
\addplot[Mapcolor,mark=*,mark options={fill=white,mark size=1.5pt}, nodes near coords={0,99}]
coordinates {(6, -4)};
\addplot[Mapcolor,mark=*,mark options={fill=white,mark size=1.5pt}, nodes near coords={1,32}]
coordinates {(8, -4)};
\addplot[Mapcolor,mark=*,mark options={fill=white,mark size=1.5pt}, nodes near coords={1,65}]
coordinates {(10, -4)};
\end{axis}
\end{tikzpicture}}	
\choice
{$2,10\ cm/s$}
{\True $4,20\ cm/s$}
{$3,70\ cm/s$}
{$4,50\ cm/s$}	
\loigiai{ }
\end{ex}
% Câu hỏi 1 — vận tốc cực đại (#TN)
\begin{ex}
Xác định vận tốc cực đại của vật dao động điều hoà có phương trình $x = 4\cos\left(5\pi t - \dfrac{\pi}{3}\right)\ (cm)$  	
\choice
{$10\pi\ cm/s$}
{$15\pi\ cm/s$}
{\True $20\pi\ cm/s$}
{$25\pi\ cm/s$}
\loigiai{
}
\end{ex}

\begin{ex}
Xác định gia tốc của vật khi đi qua vị trí có li độ $x = 1\ cm$ đối với dao động $x = 4\cos\left(5\pi t - \dfrac{\pi}{3}\right)\ (cm)$  
\choice
{$25\pi^{2}\ cm/s^{2}$}
{$-12,5\pi^{2}\ cm/s^{2}$}
{\True $-25\pi^{2}\ cm/s^{2}$}
{$-50\pi^{2}\ cm/s^{2}$}	
\loigiai{
}
\end{ex}
\begin{ex}
Một vật nhỏ dao động điều hoà có gia tốc cực đại $20\pi^{2}\ cm/s^{2}$ và tốc độ cực đại $10\pi\ cm/s$.  
Quãng đường vật đi được từ $t_1 = 1\ s$ đến $t_2 = 15,5\ s$ là
\choice
{$2,5\ m$}
{\True $2,9\ m$}
{$3,1\ m$}
{$2,7\ m$}
\loigiai{
$\omega = \dfrac{a_{\max}}{v_{\max}} = 2\pi\ rad/s$ \\
$T = \dfrac{2\pi}{\omega}=1\ s,\quad A=\dfrac{v_{\max}}{\omega}=5\ cm$ \\
$14,5\ s = 14T + 0,5T\;\Rightarrow\; S = 14\cdot4A + 2A = 2,9\ m$ \\
Vậy quãng đường cần tìm là $S = 2,9\ m$
}
\end{ex}

\setcounter{ex}{0}
\setcounter{paragraph}{0}
\PhanII
\begin{ex}
Một chất điểm dao động điều hoà có phương trình $x=4 \cos (\pi t+\dfrac{\pi}{4})\ cm$.
\choiceTF[t]
{Chu kì dao động là $4 \ s$}
{Chiều dài quỹ đạo là $4 \ cm$}
{\True Lúc $t=0$, chất điểm chuyển động theo chiều âm}
{Tốc độ của chất điểm khi qua vị trí cân bằng là $4 \ cm/s$}
\loigiai{
\begin{itemchoice}
\itemch
\itemch
\itemch
\itemch
\end{itemchoice}
}
%\haidong{4}
\end{ex} \haidong{15}
\begin{ex}
Một chất điểm dao động có phương trình $x=6 \cos \pi t$ ($x$ tính bằng $cm;\ t$ tính bằng giây).
\choiceTF[t]
{Chu kì dao động là $0,5 \ s$}
{Tần số của dao động là $2 \ Hz$}
{\True Tốc độ cực đại của chất điểm là $6 \pi\ cm/s$}
{\True Gia tốc của chất điểm có độ lớn cực đại là $6 \pi^2 \ cm/s^2$}
\loigiai{
\begin{itemchoice}
\itemch
\itemch
\itemch
\itemch
\end{itemchoice}

}
%\haidong{4}
\end{ex} \haidong{15}
\begin{ex}
Một chất điểm dao động điều hoà trên trục $O x$ có phương trình $x=8 \cos \left(\pi t+\dfrac{\pi}{2}\right)$ ($x$ tính bằng $cm,\ t$ tính bằng $s$).
\choiceTF[t]
{Chu kì dao động là $4 \ s$}
{\True Chất điểm chuyển động trên đoạn thẳng dài $16 \ cm$}
{\True Lúc $t=0$ chất điểm chuyển động theo chiều âm của trục $O x$}
{Vận tốc của chất điểm tại vị trí cân bằng là $8 \ cm/s$}
\loigiai{
\begin{itemchoice}
\itemch
\itemch
\itemch
\itemch
\end{itemchoice}
}
%\haidong{4}
\end{ex} \haidong{15}
\begin{ex}
Một chất điểm dao động điều hoà trên trục $Ox$ có phương trình vận tốc 
$v = 2\pi \cos(\pi t)\ (cm/s)$, 
trong đó gốc toạ độ trùng với vị trí cân bằng.
\choiceTF[t]
{Ở vị trí cân bằng, vận tốc của chất điểm bằng $0$}
{Pha ban đầu của dao động là $\varphi = 0$}
{\True Tại thời điểm $t=3\ s$ chất điểm đang đi theo chiều âm}
{Tại thời điểm $t = \dfrac{1}{2}\ s$, li độ của chất điểm là $x = 2\ cm$}
\loigiai{
\begin{itemchoice}
\itemch 
\itemch 
\itemch 
\itemch 
\end{itemchoice}
}
\end{ex}



\setcounter{ex}{0}
\PhanIII
\begin{ex}
Một chất điểm dao động điều hoà với phương trình $x=8 \cos 5 t(x$ tính bằng $cm, t$ tính bằng $s$). Tốc độ của chất điểm khi đi qua vị trí cân bằng bằng bao nhiêu $cm/s$?
\shortans[oly]{40}
\loigiai{
$40 \ cm/s$
}
%\haidong{4}
\end{ex} \haidong{15}
\begin{ex}
Một chất điểm dao động điều hòa dọc theo trục $O x$ trên quỹ đạo dài $10 \ cm$. Biết chất điểm thực hiện được 20 dao động toàn phần trong $5 \ s$. Tốc độ cực đại của vật trong quá trình dao động bằng bao nhiêu $m/s$ (làm tròn kết quả đến chữ số hàng phần trăm)?
\shortans[oly]{1,26}
\loigiai{
$1,26 \ m/s$
}
%\haidong{4}
\end{ex} \haidong{15}
\begin{ex}
Một vật dao động điều hoà trên trục $O x$ có vận tốc cực đại bằng $20 \ cm/s$ và gia tốc cực đại bằng $40\ cm/s^2$. Tốc độ góc của vật bằng bao nhiêu rad/s?
\shortans[oly]{2}
\loigiai{
$2\ rad/s$
}
%\haidong{4}
\end{ex} \haidong{15}
\begin{ex}
Một vật dao động điều hoà với chu kì $0,4\ s$ với chiều dài quỹ đạo là 9 cm. Lấy $\pi^2=10$. Gia tốc cực đại bằng bao nhiêu $cm/s^2$?
\shortans[oly]{1125}
\loigiai{
$1125\ cm/s^2$
}
%\haidong{5}
\end{ex} \haidong{15}
\begin{ex}
Một vật dao động điều hoà với gia tốc cực đại bằng $86,4 \ m/s^2$, vận tốc cực đại bằng $2,16 \ m/s$. Quỹ đạo chuyển động của vật là một đoạn thẳng dài có độ lớn bao nhiêu $cm$?
\shortans[oly]{10,8}
\loigiai{
$10,8 \ cm$
}
%\haidong{5}
\end{ex} \haidong{15}
\begin{ex}
Một chất điểm dao động điều hoà với phương trình li độ $x=2 \cos \left(\pi t+\dfrac{\pi}{6}\right)\ cm$ ($t$ tính bằng giây). Lấy $\pi^2=10$. Tốc độ của chất điểm ở thời điểm $t=0,5 \ s$ bằng bao nhiêu $cm/s$ (làm tròn kết quả đến chữ số hàng phần trăm)?
\shortans[oly]{5,44}
\loigiai{
$5,44 \ cm/s$
}
%\haidong{5}
\end{ex} \haidong{15}